\chapter{Fungsi Acuan}\label{ch:g10:reffunc}

Segelintir \omterm{def:g10:functions:function}{fungsi} sederhana --- afin, pangkat dua, kebalikan, akar kuadrat,
pangkat tiga --- muncul di mana-mana, sendirian atau berpadu. Hafal di luar
kepala \omterm{def:g10:functions:graph}{grafik} dan variasinya mengubah banyak soal menjadi sebuah sketsa
singkat. Bab ini menelaah masing-masing secara berurutan dan membuktikan
variasinya dengan teknik pembandingan dari \cref{ch:g10:functions}.

\section{Fungsi afin}

\begin{definition}[Fungsi afin]\label{def:g10:reffunc:affine}
Sebuah \emph{fungsi afin}\index{fungsi afin} adalah \omterm{def:g10:functions:function}{fungsi} berbentuk
\[
f(x) = mx + p,
\]
dengan $m$ dan $p$ \omterm{def:g10:numbers:sets}{bilangan real} tetap. \omterm{def:g10:functions:graph}{Grafiknya} berupa garis lurus:
$m$ adalah \emph{gradien}\index{gradien} dan $p$ adalah
\emph{titik potong sumbu-$y$}\index{titik potong sumbu-$y$} (garis itu
memotong sumbu tegak di $(0, p)$). Bila $p = 0$ \omterm{def:g10:functions:function}{fungsi} itu disebut
\emph{linear}: $f(x) = mx$ menyatakan kesebandingan.
\end{definition}

\begin{proposition}[Gradien dan variasi]\label{prop:g10:reffunc:affine}
Misalkan $f(x) = mx + p$.
\begin{enumerate}
  \item Untuk sebarang dua masukan berbeda $u \neq v$:
        $m = \dfrac{f(v) - f(u)}{v - u}$ (\omterm{def:g10:reffunc:affine}{gradien} adalah perubahan
        keluaran per satu satuan perubahan masukan).
  \item Jika $m > 0$, maka $f$ tegas naik pada $\R$; jika $m < 0$, tegas
        turun; jika $m = 0$, tetap.
\end{enumerate}
\end{proposition}

\begin{proof}
1. Hitunglah $f(v) - f(u) = (mv + p) - (mu + p) = m(v - u)$, lalu bagilah
dengan $v - u \neq 0$.

2. Jika $u < v$, maka $f(v) - f(u) = m(v-u)$ bertanda sama dengan $m$, sebab
$v - u > 0$: untuk $m$ positif diperoleh $f(u) < f(v)$ (naik), untuk $m$
negatif diperoleh $f(u) > f(v)$ (turun).
\end{proof}

\begin{omfigure}
\begin{tikzpicture}
\begin{axis}[omaxis,
  width=7.8cm, height=5.6cm,
  xmin=-3.4, xmax=3.4, ymin=-2.8, ymax=4.4,
  xtick={-2,-1,1,2}, ytick={-1,1,2,3}]
  \addplot[omDef, thick, domain=-3:2.6] {0.5*x + 1};
  \addplot[omThm, thick, domain=-2.2:3] {-x + 1};
  \node[omDef, font=\small] at (axis cs:2.5,2.7) {$y = \tfrac12 x + 1$};
  \node[omThm, font=\small] at (axis cs:-2.0,2.6) {$y = -x + 1$};
  \draw[dashed] (axis cs:1,1.5) -- (axis cs:2,1.5) -- (axis cs:2,2);
  \node[font=\footnotesize, below] at (axis cs:1.5,1.5) {$1$};
  \node[font=\footnotesize, right] at (axis cs:2,1.75) {$m$};
\end{axis}
\end{tikzpicture}

{\small Dua \omterm{def:g10:reffunc:affine}{fungsi afin}: naik untuk $m = \frac12 > 0$, turun untuk
$m = -1 < 0$. Bergerak satu satuan ke kanan mengubah $y$ sebesar $m$.}
\end{omfigure}

\begin{example}\label{ex:g10:reffunc:affine}
Carilah \omterm{def:g10:reffunc:affine}{fungsi afin} yang \omterm{def:g10:functions:graph}{grafiknya} melalui $A(1, 3)$ dan
$B(4, 9)$. \omterm{def:g10:reffunc:affine}{Gradiennya} adalah
\[
m = \frac{9 - 3}{4 - 1} = \frac{6}{3} = 2 .
\]
Jadi $f(x) = 2x + p$, dan $f(1) = 3$ memberi $2 + p = 3$, sehingga $p = 1$:
$f(x) = 2x + 1$. Periksa dengan $B$: $f(4) = 9$.
\end{example}

\section{Fungsi pangkat dua}

\begin{proposition}[Fungsi pangkat dua]\label{prop:g10:reffunc:square}
\omterm{def:g10:functions:function}{Fungsi} $f(x) = x^2$, terdefinisi pada $\R$:
\begin{enumerate}
  \item tegas turun pada $\intoc{-\infty}{0}$ dan tegas
        naik pada $\intco{0}{+\infty}$, dengan \omterm{def:g10:functions:extrema}{minimum} $0$ di $x = 0$;
  \item memenuhi $f(-x) = f(x)$ untuk semua $x$: \omterm{def:g10:functions:graph}{grafiknya}, sebuah
        \emph{parabola}\index{parabola}, simetris terhadap sumbu
        tegak.
\end{enumerate}
\end{proposition}

\begin{proof}
1. Misalkan $0 \leq u < v$. Maka
\[
f(v) - f(u) = v^2 - u^2 = (v - u)(v + u) > 0,
\]
sebab $v - u > 0$ dan $v + u > 0$: jadi $f$ tegas naik pada
$\intco{0}{+\infty}$. Jika $u < v \leq 0$, maka $v - u > 0$ tetapi
$v + u < 0$, sehingga $f(v) - f(u) < 0$: tegas turun. Akhirnya
$x^2 \geq 0 = f(0)$ untuk semua $x$.

2. $(-x)^2 = x^2$; jadi titik $(x, x^2)$ dan $(-x, x^2)$, yang saling
bercermin terhadap sumbu tegak, sama-sama terletak pada \omterm{def:g10:functions:graph}{grafiknya}.
\end{proof}

\begin{remark}[Kuadrat dan urutan]
Karena \omterm{def:g10:functions:function}{fungsi} pangkat dua turun di sisi negatif, mengambil kuadrat
\emph{membalik} urutan bilangan negatif: $-3 < -2$ tetapi $9 > 4$.
Jangan pernah mengkuadratkan kedua ruas suatu pertidaksamaan tanpa memeriksa tandanya.
\end{remark}

\section{Fungsi kebalikan}

\begin{proposition}[Fungsi kebalikan]\label{prop:g10:reffunc:inverse}
\omterm{def:g10:functions:function}{Fungsi} $f(x) = \dfrac1x$, terdefinisi untuk $x \neq 0$:
\begin{enumerate}
  \item tegas turun pada $\intoo{-\infty}{0}$ dan tegas
        turun pada $\intoo{0}{+\infty}$;
  \item memenuhi $f(-x) = -f(x)$: \omterm{def:g10:functions:graph}{grafiknya}, sebuah
        \emph{hiperbola}\index{hiperbola}, simetris terhadap titik asal.
\end{enumerate}
\end{proposition}

\begin{proof}
1. Misalkan $0 < u < v$. Maka
\[
f(v) - f(u) = \frac1v - \frac1u = \frac{u - v}{uv} < 0,
\]
sebab $u - v < 0$ dan $uv > 0$: tegas turun pada
$\intoo{0}{+\infty}$. Pada $\intoo{-\infty}{0}$ pecahan yang sama mempunyai
$u - v < 0$ dan lagi-lagi $uv > 0$ (hasil kali dua bilangan negatif), jadi $f$
turun di sana juga.

2. $\frac{1}{-x} = -\frac1x$.
\end{proof}

\begin{remark}
Hati-hati: $\frac1x$ \emph{tidak} turun pada seluruh \omterm{def:g10:functions:function}{daerah asalnya}. Dari
$u = -1$ ke $v = 1$ nilainya melompat dari $-1$ naik ke $1$. Kedua cabangnya
harus ditelaah terpisah.
\end{remark}

\begin{omfigure}
\begin{tikzpicture}
\begin{axis}[omaxis,
  width=5.6cm, height=5.2cm,
  xmin=-3.2, xmax=3.2, ymin=-2.6, ymax=6.4,
  xtick={-2,-1,1,2}, ytick={1,4},
  title={$y = x^2$}]
  \addplot[omDef, thick, domain=-2.45:2.45, samples=90] {x^2};
\end{axis}
\end{tikzpicture}\hfill
\begin{tikzpicture}
\begin{axis}[omaxis,
  width=5.6cm, height=5.2cm,
  xmin=-3.4, xmax=3.4, ymin=-3.4, ymax=3.4,
  xtick={-1,1}, ytick={-1,1},
  title={$y = \frac1x$}]
  \addplot[omDef, thick, domain=0.3:3.1, samples=80] {1/x};
  \addplot[omDef, thick, domain=-3.1:-0.3, samples=80] {1/x};
\end{axis}
\end{tikzpicture}

{\small Parabola $y = x^2$ (simetris terhadap sumbu tegak) dan
hiperbola $y = \frac1x$ (dua cabang, simetris terhadap titik asal).}
\end{omfigure}

\section{Akar kuadrat dan pangkat tiga}

\begin{proposition}[Fungsi akar kuadrat]\label{prop:g10:reffunc:sqrt}
\omterm{def:g10:functions:function}{Fungsi} $f(x) = \sqrt{x}$, terdefinisi pada $\intco{0}{+\infty}$, bersifat
tegas naik.
\end{proposition}

\begin{proof}
Misalkan $0 \leq u < v$. Kalikan dan bagilah dengan bentuk \emph{sekawan}:
\[
\sqrt v - \sqrt u
= \frac{(\sqrt v - \sqrt u)(\sqrt v + \sqrt u)}{\sqrt v + \sqrt u}
= \frac{v - u}{\sqrt v + \sqrt u} > 0,
\]
sebab $v - u > 0$ dan $\sqrt v + \sqrt u > 0$ (perhatikan $v > 0$, jadi
$\sqrt v > 0$).
\end{proof}

\begin{proposition}[Fungsi pangkat tiga]\label{prop:g10:reffunc:cube}
\omterm{def:g10:functions:function}{Fungsi} $f(x) = x^3$, terdefinisi pada $\R$, bersifat tegas naik, dan
\omterm{def:g10:functions:graph}{grafiknya} simetris terhadap titik asal.
\end{proposition}
\admitted

\begin{omfigure}
\begin{tikzpicture}
\begin{axis}[omaxis,
  width=5.6cm, height=5.0cm,
  xmin=-0.7, xmax=5.4, ymin=-0.7, ymax=2.9,
  xtick={1,4}, ytick={1,2},
  title={$y = \sqrt x$}]
  \addplot[omDef, thick, domain=0:5.1, samples=100] {sqrt(x)};
\end{axis}
\end{tikzpicture}\hfill
\begin{tikzpicture}
\begin{axis}[omaxis,
  width=5.6cm, height=5.0cm,
  xmin=-2.2, xmax=2.2, ymin=-3.4, ymax=3.4,
  xtick={-1,1}, ytick={-1,1},
  title={$y = x^3$}]
  \addplot[omDef, thick, domain=-1.48:1.48, samples=90] {x^3};
\end{axis}
\end{tikzpicture}

{\small Akar kuadrat (terdefinisi hanya untuk $x \geq 0$) dan pangkat tiga,
keduanya naik.}
\end{omfigure}

\begin{proposition}[Membandingkan $x$, $x^2$ dan $\sqrt x$]\label{prop:g10:reffunc:compare}
Untuk $0 < x < 1$:\ \ $x^2 < x < \sqrt x$.\quad
Untuk $x > 1$:\ \ $\sqrt x < x < x^2$.\quad
Di $x = 0$ dan $x = 1$ ketiga nilainya sama.
\end{proposition}

\begin{proof}
Andaikan $0 < x < 1$. Mengalikan $x < 1$ dengan $x > 0$ memberi $x^2 < x$.
Selanjutnya, $x < \sqrt x$: karena kedua ruas positif dan pengkuadratan
bersifat naik pada bilangan positif, ketaksamaan ini setara dengan
$x^2 < x$, yang baru saja dibuktikan. Untuk $x > 1$, mengalikan $x > 1$ dengan
$x$ memberi $x^2 > x$, dan $\sqrt x < x$ menyusul dengan alasan pengkuadratan yang sama.
\end{proof}

\begin{omfigure}
\begin{tikzpicture}
\begin{axis}[omaxis,
  width=7.4cm, height=5.8cm,
  xmin=-0.35, xmax=2.5, ymin=-0.35, ymax=2.9,
  xtick={1,2}, ytick={1,2}]
  \addplot[omThm, thick, domain=0:1.65, samples=80] {x^2};
  \addplot[omDef, thick, domain=0:2.35] {x};
  \addplot[omMeth, thick, domain=0:2.35, samples=80] {sqrt(x)};
  \fill (axis cs:1,1) circle (1.5pt);
  \node[omThm, font=\small, right] at (axis cs:1.45,2.4) {$y = x^2$};
  \node[omDef, font=\small, right] at (axis cs:2.0,1.85) {$y = x$};
  \node[omMeth, font=\small, right] at (axis cs:2.0,1.28) {$y = \sqrt x$};
\end{axis}
\end{tikzpicture}

{\small Ketiga kurva berpotongan di $(0,0)$ dan $(1,1)$ lalu bertukar urutan
di sana: antara $0$ dan $1$, $x^2$ yang terkecil dan $\sqrt x$ yang
terbesar; setelah $1$ urutannya terbalik.}
\end{omfigure}

\begin{method}[Membandingkan peta]\label{met:g10:reffunc:compare}
Untuk membandingkan $f(a)$ dan $f(b)$ pada sebuah \omterm{def:g10:functions:function}{fungsi} acuan $f$:
\begin{enumerate}
  \item tempatkan $a$ dan $b$ pada \omterm{def:g10:numbers:interval}{interval} yang variasi $f$-nya sudah
        diketahui;
  \item jika $f$ naik di sana, \omterm{def:g10:functions:function}{petanya} berurutan sama dengan
        masukannya; jika $f$ turun, urutannya terbalik;
  \item jika $a$ dan $b$ tidak berada pada \omterm{def:g10:numbers:interval}{interval} kemonotonan yang sama,
        uruslah tandanya lebih dahulu (misalnya untuk pangkat dua: satu
        masukan negatif dan satu masukan positif).
\end{enumerate}
\end{method}

\begin{example}
Bandingkan $\dfrac{1}{2.3}$ dan $\dfrac{1}{2.4}$: kedua masukannya berada di
$\intoo{0}{+\infty}$, tempat \omterm{def:g10:functions:function}{fungsi} kebalikan turun, dan
$2.3 < 2.4$, jadi $\dfrac{1}{2.3} > \dfrac{1}{2.4}$. Bandingkan $(-1.2)^2$
dan $(-1.5)^2$: pada $\intoc{-\infty}{0}$ pangkat dua turun, dan
$-1.5 < -1.2$, jadi $(-1.5)^2 > (-1.2)^2$ --- memang $2.25 > 1.44$.
\end{example}

\section{Latihan}

\begin{exercise}[$\star$]\label{exo:g10:reffunc:1}
Untuk setiap \omterm{def:g10:reffunc:affine}{fungsi afin} berikut, berikan \omterm{def:g10:reffunc:affine}{gradiennya}, titik potong
sumbu-$y$-nya, dan sebutkan apakah \omterm{def:g10:functions:function}{fungsi} itu naik atau turun:
\[
f(x) = 3x - 2, \qquad
g(x) = -\tfrac12 x + 4, \qquad
h(x) = 7, \qquad
k(x) = -x .
\]
\end{exercise}

\begin{exercise}[$\star$]\label{exo:g10:reffunc:2}
Carilah \omterm{def:g10:reffunc:affine}{fungsi afin} yang \omterm{def:g10:functions:graph}{grafiknya} melalui $(2, 1)$ dan
$(5, 10)$; lalu yang melalui $(-1, 4)$ dan $(3, -4)$.
\end{exercise}

\begin{exercise}[$\star$]\label{exo:g10:reffunc:3}
Tanpa kalkulator, bandingkan:
\[
(3.1)^2 \text{ dan } (3.2)^2; \qquad
(-2.7)^2 \text{ dan } (-2.8)^2; \qquad
\frac{1}{5.1} \text{ dan } \frac{1}{5.2}; \qquad
\sqrt{17} \text{ dan } \sqrt{15}.
\]
\end{exercise}

\begin{exercise}[$\star$]\label{exo:g10:reffunc:4}
Dengan memakai \omterm{def:g10:functions:graph}{grafik} \omterm{def:g10:functions:function}{fungsi} pangkat dua, selesaikan $x^2 = 16$, lalu
$x^2 \leq 16$, lalu $x^2 > 9$.
\end{exercise}

\begin{exercise}[$\star$]\label{exo:g10:reffunc:5}
Misalkan $x \in \intoo{0}{1}$. Urutkan bilangan $x$, $x^2$, $x^3$ dan
$\sqrt x$ dari yang terkecil sampai yang terbesar, lalu periksa jawabanmu
dengan $x = 0.25$.
\end{exercise}

\begin{exercise}[$\star\star$]\label{exo:g10:reffunc:6}
Selesaikan secara grafis, lalu secara aljabar: $\dfrac1x = 2$, dan
$\dfrac1x < 2$ untuk $x > 0$. Apa yang berubah bila $x < 0$ juga diizinkan?
\end{exercise}

\begin{exercise}[$\star\star$]\label{exo:g10:reffunc:7}
Dengan mengetahui bahwa \omterm{def:g10:functions:function}{fungsi} pangkat dua turun pada $\intoc{-\infty}{0}$ dan
naik pada $\intco{0}{+\infty}$, apitlah $x^2$ bila:
\[
\text{(a) } x \in \intcc{2}{5};
\qquad
\text{(b) } x \in \intcc{-3}{-1};
\qquad
\text{(c) } x \in \intcc{-2}{3}.
\]
(Pada kasus (c), berhati-hatilah: \omterm{def:g10:functions:extrema}{minimum} $x^2$ tidak berada di titik ujung.)
\end{exercise}

\begin{exercise}[$\star\star$]\label{exo:g10:reffunc:8}
Sebuah perusahaan taksi menarik biaya tetap $4$ ditambah $1.5$ per kilometer;
perusahaan lain tanpa biaya tetap dan $2$ per kilometer. Modelkan masing-masing
harga dengan \omterm{def:g10:reffunc:affine}{fungsi afin} dari jarak $x$, gambarlah kedua garisnya, lalu carilah
mulai jarak berapa perusahaan pertama lebih murah.
\end{exercise}

\begin{exercise}[$\star\star$]\label{exo:g10:reffunc:9}
Tunjukkan bahwa untuk semua $a, b \geq 0$: $\sqrt{a + b} \leq \sqrt a + \sqrt b$.
(Petunjuk: kedua ruas taknegatif, jadi bandingkan kuadratnya.) Kapan
kesamaannya berlaku?
\end{exercise}

\begin{exercise}[$\star\star\star$]\label{exo:g10:reffunc:10}
Misalkan $f(x) = \dfrac{2x + 1}{x - 1}$, terdefinisi untuk $x \neq 1$.
\begin{enumerate}
  \item Tunjukkan bahwa $f(x) = 2 + \dfrac{3}{x - 1}$ untuk semua $x \neq 1$.
  \item Simpulkan variasi $f$ pada $\intoo{1}{+\infty}$ dari variasi
        \omterm{def:g10:functions:function}{fungsi} kebalikan.
\end{enumerate}
\end{exercise}

\section{Soal: Hukum alam berbicara dengan fungsi acuan}

\begin{problem}[{Soal akhir pekan --- jarak pengereman tumbuh seperti
$v^2$, tuas menuruti $1/d$, planet mengikuti $a^{3/2}$: membaca
fisika dengan \omterm{def:g10:functions:function}{fungsi} acuan}]
\label{pb:g10:reffunc:1}
Bukalah sebuah buku fisika dan \omterm{def:g10:functions:graph}{grafik} yang itu-itu juga muncul di setiap
halaman: parabola, hiperbola, kurva akar. Alam menuliskan
hukumnya dengan \omterm{def:g10:functions:function}{fungsi} acuan bab ini --- dan mengenali
bentuknya (\cref{prop:g10:reffunc:square},
\cref{prop:g10:reffunc:inverse}, \cref{prop:g10:reffunc:sqrt},
\cref{prop:g10:reffunc:cube}) sudah cukup untuk mengerem mobil, menyeimbangkan
tuas, dan mengukur waktu edar planet. Puncaknya adalah hukum ketiga
Kepler, yang dibaca langsung dari data tata surya.

\medskip
\noindent\textbf{Bagian I --- Hukum kuadrat pengereman.}
Sebuah aturan praktis untuk jarak pengereman mobil di jalan kering:
$d = \left(\frac{v}{10}\right)^{2}$ meter, pada laju $v$ km/jam.
\begin{enumerate}
  \item Hitunglah jarak pengereman pada $30$, $50$, $90$ dan
        $130$ km/jam.
  \item Melipatduakan lajunya mengalikan jarak pengereman dengan
        berapa? Jelaskan dari sifat penskalaan \omterm{def:g10:functions:function}{fungsi} pangkat dua,
        lalu periksa pada $50 \to 100$ km/jam.
  \item Forensik: bekas rem terukur $50$ m. Laju berapa yang
        dituduhkan rumus itu kepada pengemudinya (sampai satuan km/jam)? \omterm{def:g10:functions:function}{Fungsi}
        acuan mana yang menjawab --- dan pada \omterm{def:g10:functions:function}{daerah asal} mana
        pembacaannya tidak bermakna ganda
        (\cref{prop:g10:reffunc:sqrt})?
  \item Bandingkan tambahan jarak yang disebabkan oleh \emph{satu} km/jam
        lebih cepat, pada laju rendah dan pada laju tinggi: hitunglah
        $d(31) - d(30)$ dan $d(91) - d(90)$. Sifat parabola yang mana
        (\cref{prop:g10:reffunc:square}) yang digambarkan oleh kedua
        jawaban itu?
  \item Pengereman yang sesungguhnya menambahkan waktu reaksi --- sekitar satu detik,
        yang selama itu mobil menempuh $0.28v$ meter:
        $D(v) = 0.28v + \left(\frac{v}{10}\right)^2$. Hitunglah
        $D(50)$ dan $D(130)$. Manakah dari kedua suku itu --- yang afin
        atau yang pangkat dua --- yang menguasai laju dalam kota, dan mana yang
        menguasai laju di jalan tol?
\end{enumerate}

\medskip
\noindent\textbf{Bagian II --- Hiperbola dalam kehidupan sehari-hari.}
\begin{enumerate}[resume]
  \item Sebuah tuas seimbang apabila gaya $\times$ jarak bernilai
        sama pada kedua sisinya. Seorang anak $60$ kg duduk $2$ m dari
        poros; gaya penyeimbang pada jarak $d$ di sisi lain
        adalah $F(d) = \frac{120}{d}$ (dalam satuan setara kilogram).
        Hitunglah $F(0.5)$, $F(1)$, $F(3)$.
  \item Apa yang dikatakan variasi \omterm{def:g10:functions:function}{fungsi} kebalikan
        (\cref{prop:g10:reffunc:inverse}) tentang tuas ---
        dan mengapa kesesumbaran Archimedes (``beri aku tempat
        berpijak dan akan kugerakkan Bumi'') bersembunyi di
        ekor hiperbola? Apa yang melarang $d = 0$?
  \item Bunyi dan cahaya meredup mengikuti \emph{kuadrat}
        jaraknya: pada $1$ m dari sebuah lampu intensitasnya $100$
        satuan; berikan nilainya pada $2$, $3$ dan $10$ m. Jelaskan
        pangkat $2$ itu dengan sebuah bola: pada luas berapa cahaya
        itu tersebar pada jarak $d$ (soal batu nisan Archimedes di jilid sebelumnya)?
  \item Sebuah bus untuk karyawisata berbiaya $600$ euro, dibagi
        rata di antara $n$ peserta. Buatlah tabel biaya per
        kepala untuk $n = 10, 20, 30$, carilah berapa peserta yang
        menurunkannya menjadi $25$ euro atau kurang, lalu nyatakan apa yang
        dijanjikan --- dan ditolak --- oleh bentuk hiperbola ketika $n$
        membesar.
  \item Mencetak $x$ poster berbiaya $2x + 3$ euro ($3$ euro untuk
        penyiapan). Tunjukkan bahwa biaya \emph{rata-rata} per poster adalah
        $a(x) = 2 + \frac3x$, uraikan variasinya, lalu
        tafsirkan asimtot mendatarnya secara ekonomi.
        (Bandingkan dengan aljabar \cref{exo:g10:reffunc:10}.)
\end{enumerate}

\medskip
\noindent\textbf{Bagian III --- Keselarasan Kepler.}
Untuk setiap planet, misalkan $a$ jarak rata-ratanya ke Matahari (dalam
satuan astronomi, Bumi $= 1$) dan $T$ kala edarnya (dalam tahun).
Data: Mars $a = 1.52$, $T = 1.88$; Jupiter $a = 5.20$,
$T = 11.86$; Saturnus $a = 9.54$, $T = 29.4$.
\begin{enumerate}[resume]
  \item Hitunglah $a^3$ dan $T^2$ untuk Bumi, Mars dan Jupiter.
        Apa yang diperhatikan Kepler pada tahun 1618?
  \item Nyatakan hukum itu sebagai \omterm{def:g10:functions:function}{fungsi}: $T = \sqrt{a^3} = a\sqrt a$.
        Ujilah pada Saturnus.
  \item Dua ramalan: sebuah asteroid mengorbit pada $a = 4$ satuan astronomi ---
        carilah kala edarnya (bilangannya keluar bulat); sebuah komet
        kembali setiap $27$ tahun --- carilah jarak rata-ratanya
        (carilah sebuah pangkat tiga sempurna).
  \item Tiga \omterm{def:g10:functions:function}{fungsi} acuan yang mana yang dijalin bersama oleh hukum itu?
        Lalu aturan penskalaannya: ketika $a$ dikalikan
        $4$, apa yang terjadi pada $T$? (Periksa dengan asteroid tadi
        terhadap Bumi.)
  \item Kepler membaca hukum ini dalam data Tycho Brahe tujuh puluh tahun
        sebelum gravitasi Newton menjelaskannya. Dalam satu
        kalimat: apa yang dikatakan peristiwa ini tentang daya
        mengenali sebuah \omterm{def:g10:functions:function}{fungsi} acuan di dalam tabel bilangan?
\end{enumerate}

\medskip
\noindent\textbf{Bagian IV --- Kura-kura dan kelinci.}
\begin{enumerate}[resume]
  \item Urutkan $\sqrt x$, $x$ dan $x^2$ pada $\intoo{0}{1}$ dan pada
        $\intoo{1}{+\infty}$
        (\cref{prop:g10:reffunc:compare}), lalu periksa kedua
        urutannya di $x = 0.25$ dan $x = 4$.
  \item Selesaikan $\sqrt x > x$ selengkapnya (untuk $x \geq 0$, kuadratkan
        dengan hati-hati lalu tuntaskan dengan \omterm{def:g10:algebra:signtable}{tabel tanda}).
  \item Masukkan pangkat tiga: urutkan $x$, $x^2$, $x^3$ di
        $x = 0.5$ dan di $x = 2$, lalu carilah semua titik tempat
        \omterm{def:g10:functions:function}{fungsi} pangkat dua dan pangkat tiga berpotongan.
  \item Dua skema tabungan membayar, setelah $t$ tahun, $100\sqrt t$
        euro (skema A) atau $10 t^2$ euro (skema B). Tunjukkan bahwa
        B menyalip A ketika $t^3 = 100$, lalu berikan waktu
        penyalipannya sampai satuan bulan. Apa moralnya tentang akar melawan
        pangkat dalam jangka panjang?
  \item Penutup: untuk setiap hukum dalam soal ini --- pengereman
        ($v^2$), tuas ($1/d$), lampu ($1/d^2$), Kepler
        ($a^{3/2}$), skema A ($\sqrt t$) --- sebutkan dalam satu anak kalimat
        ciri \omterm{def:g10:functions:graph}{grafik} acuannya yang mengusung
        makna fisisnya (kenaikan yang makin curam, asimtot tegak,
        bola yang menyebar, pangkat penskalaan, pertumbuhan yang mendatar).
        Kesimpulannya: \omterm{def:g10:functions:function}{fungsi} acuan adalah abjadnya;
        alam menulis dengan abjad itu.

\end{enumerate}
\end{problem}
